\subsection{Минимальная степень слайса: точный результат, независимый от выпрямления}

Рассмотрим над произвольным полем $k$ характеристики нуль дифференцирование
\[
 D=t\partial_u+(3u^2-3)\partial_x+(4u^3-8u)\partial_y
 +(5u^4-10)\partial_z
 \quad\text{на }k[t,u,x,y,z].
\]
Слайсом называется полином $s$ такой, что $D(s)=1$.
Полная степень ниже является обычной полной степенью по всем пяти переменным
$t,u,x,y,z$, включая $t$.

\paragraph{Теорема.}
Минимальная полная степень полиномиального слайса для $D$ равна $7$.
Утверждение не использует выпрямитель из репозитория. Здесь не заявляется,
что этот результат отсутствует в литературе: установлена его самостоятельная
точная проверка с проверяемым сертификатом нижней границы.

\paragraph{Верхняя граница.}
Положим $f=u^3-3u$, $g=u^4-4u^2$, $h=u^5-10u$. Тогда
\begin{align*}
 \sigma={}&(g+2)z+(h-5f)y-(3f^2+5g+7)x\\
 &+t(3fx^2+5xy-yz)-t^2x^3.
\end{align*}
Иными словами,
\begin{align*}
 \sigma={}&-(3u^6-13u^4+7u^2+7)x
 +(u^5-5u^3+5u)y+(u^4-4u^2+2)z\\
 &+t\bigl(3(u^3-3u)x^2+5xy-yz\bigr)-t^2x^3.
\end{align*}
У этого полинома $15$ ненулевых мономов, полная степень $7$,
и непосредственное дифференцирование даёт $D(\sigma)=1$.
Полезное происхождение формулы: если
\[
 F_0=yz-x^3-5xy+2z-7x,\qquad
 V=(f,g,h)-t(x,y,z),
\]
то $F_0(f(u),g(u),h(u))=u$ и $\sigma=(u-F_0(V))/t$.
Последнее выражение является точным полиномиальным частным, а не
рациональной функцией с неустранённым полюсом.

\paragraph{Элементарное доказательство нижней границы.}
Введём вес: $\operatorname{wt}(t)=1$, $\operatorname{wt}(u)=0$,
$\operatorname{wt}(x)=\operatorname{wt}(y)=\operatorname{wt}(z)=-1$.
Дифференцирование $D$ повышает вес на единицу. Поэтому если бы существовал
слайс полной степени не выше $6$, его компонента веса $-1$ также была бы
слайсом той же степени. Обозначим $X=(x,y,z)$ и
$D_0=f'(u)\partial_x+g'(u)\partial_y+h'(u)\partial_z$.
Компонента веса $-1$ имеет вид
\[
 s=s_0(u,X)+t s_1(u,X)+t^2s_2(u,X),
\]
где $s_i$ однороден по $X$ степени $i+1$, а
$\deg_u s_i\leq 5-2i$. Сравнение коэффициентов при степенях $t$ даёт
\[
 D_0s_0=1,\qquad \partial_us_0+D_0s_1=0,\qquad
 \partial_us_1+D_0s_2=0,\qquad \partial_us_2=0.
\]
Последнее равенство означает, что $s_2$ --- кубическая форма от $X$
с постоянными коэффициентами. В третьем равенстве степень
$\partial_us_1$ по $u$ не выше $2$. Коэффициенты при $u^4$ и $u^3$
поэтому дают $\partial_zs_2=\partial_ys_2=0$, то есть $s_2=cx^3$.
Интегрируя третье равенство, получаем
\[
 s_1=-3c(u^3-3u)x^2+K_2(X),
\]
где $K_2$ --- постоянная квадратичная форма. Коэффициент при $u^5$
во втором равенстве равен $-18cx$, поскольку $\deg_u\partial_us_0\leq4$
и $\deg_u D_0K_2\leq4$. Значит, $c=0$. Следовательно,
\[
 s_1=K_2(X),\qquad
 s_0=L_1(X)-\gamma(u)\cdot\nabla K_2(X)
\]
для некоторой линейной формы $L_1$. Первое равенство теперь принимает вид
\[
 \frac{d}{du}\left(L_1(\gamma(u))-K_2(\gamma(u))\right)=1.
\]
Так как $\gamma(0)=0$, получаем
$L_1(\gamma(u))-K_2(\gamma(u))=u$.
Но нечётная часть ограничения любого полинома степени не выше $2$
на $\gamma$ лежит в линейной оболочке $f,h,fg,gh$.
Их степени $3,5,7,9$ попарно различны, поэтому никакая их линейная
комбинация не может равняться $u$. Противоречие доказывает нижнюю границу.
Доказательство работает над любым полем характеристики нуль.

\paragraph{Независимый вычислительный сертификат нижней границы.}
Для монома $m$ обозначим через $[m]A$ его коэффициент в полиноме $A$.
Определим $\mathbb Q$-линейный функционал
\[
 L(A)=\sum_{m\in\mathcal S} c_m[m]A,
\]
где все $31$ ненулевой вес $c_m$ приведены в таблице ниже;
веса остальных мономов равны нулю.
Точная проверка по правилу Лейбница устанавливает
\[
 L(1)=1,\qquad L(Dm)=0
 \quad\text{для каждого монома }m\text{ полной степени }\leq 6.
\]
Таких мономов ровно $\binom{11}{5}=462$.
Для любого полинома $s$ степени не выше $6$ отсюда следует $L(Ds)=0$.
Следовательно, равенство $Ds=1$ невозможно.
Это вместе с явным $\sigma$ доказывает теорему.

\begin{center}
\renewcommand{\arraystretch}{1.25}
\begin{tabular}{r r | r r}
$m$ & $c_m$ & $m$ & $c_m$ \\
\hline
$1$ & $1$ & $t^{2}z^{2}$ & $\frac{13375}{126}$ \\
$u^{2}$ & $1$ & $t^{2}y^{2}$ & $\frac{128}{63}$ \\
$u^{4}$ & $2$ & $t^{2}xz$ & $\frac{725}{42}$ \\
$u^{6}$ & $\frac{34}{7}$ & $t^{2}x^{2}$ & $\frac{27}{7}$ \\
$u^{8}$ & $\frac{116}{9}$ & $t^{2}uyz$ & $\frac{50}{9}$ \\
$tuz$ & $-\frac{50}{7}$ & $t^{2}uxy$ & $\frac{5}{3}$ \\
$tux$ & $-\frac{3}{2}$ & $t^{2}u^{2}z^{2}$ & $\frac{3875}{42}$ \\
$tu^{2}y$ & $-\frac{8}{7}$ & $t^{2}u^{2}y^{2}$ & $-\frac{152}{21}$ \\
$tu^{3}z$ & $-\frac{100}{9}$ & $t^{2}u^{2}xz$ & $\frac{995}{42}$ \\
$tu^{3}x$ & $-\frac{15}{7}$ & $t^{2}u^{2}x^{2}$ & $\frac{27}{7}$ \\
$tu^{4}y$ & $-\frac{160}{63}$ & $t^{2}u^{3}yz$ & $\frac{1430}{63}$ \\
$tu^{5}z$ & $-\frac{6275}{252}$ & $t^{2}u^{3}xy$ & $\frac{143}{21}$ \\
$tu^{5}x$ & $-\frac{57}{14}$ & $t^{2}u^{4}z^{2}$ & $\frac{13375}{63}$ \\
$tu^{6}y$ & $-\frac{149}{63}$ & $t^{2}u^{4}y^{2}$ & $-\frac{304}{21}$ \\
$tu^{7}z$ & $-\frac{12575}{252}$ & $t^{2}u^{4}xz$ & $\frac{1625}{42}$ \\
$tu^{7}x$ & $-\frac{859}{84}$ &  &  \\
\end{tabular}
\end{center}

Все веса рациональны. Продолжение $L$ по $k$-линейности доказывает ту же
нижнюю границу над любым полем $k$ характеристики нуль, в частности
над $\mathbb Q(\sqrt{-3})$ и над $\mathbb C$.
Нижняя граница относится ко всем полиномиальным слайсам, а не только
к выбранному анзацу или к полиномам, полученным из $F_0$.

\paragraph{Воспроизводимость.}
Файл \texttt{verify\_minimal\_slice.py} использует только стандартную
библиотеку Python и точные рациональные числа. Он читает
\texttt{no\_slice\_degree6\_certificate.json}, проверяет все $462$ равенства,
нормировку $L(1)=1$ и непосредственно вычисляет $D(\sigma)$.
Воспроизводимый сертификат не зависит от решения системы линейных
уравнений, от SymPy и от формул выпрямления. Дополнительно выполнена
независимая проверка с другим кодом правила Лейбница; объединение
мономов в образах $Dm$ имеет $812$ элементов и максимальную степень $9$.

\paragraph{Связь с явными координатами ядра.}
Если уже установлен выпрямитель $F=(P,Q,R)$ и построены
$p=P(V)$, $q=Q(V)/t$, $r=R(V)/t$, $s=(u-p)/t$, то компактный слайс
$\sigma$ можно использовать вместо $s$:
\[
 \sigma-s=\frac{p-F_0(F^{-1}(p,tq,tr))}{t}\in\ker D.
\]
Это полиномиальное выражение в $t,p,q,r$, поэтому
$(t,p,q,r,\sigma)$ также является системой полиномиальных координат.
В отличие от утверждения о минимальной степени, этот последний вывод
использует проверенный выпрямитель.
